\documentclass[11pt,a4paper]{article}
\usepackage[T1]{fontenc}
\usepackage[utf8]{inputenc}
\usepackage[english]{babel}
\usepackage{amsmath,amssymb}
\usepackage[margin=2cm]{geometry}
\setlength{\parindent}{0pt}
\setlength{\emergencystretch}{3em}
\allowdisplaybreaks

\title{About \texttt{book\_exponentials\_logarithms.tex}}
\author{}
\date{}
\begin{document}
\maketitle
This note explains what the test data in \texttt{book\_exponentials\_logarithms.tex} is used for. That file is input for the engine in \texttt{main.c}; it is read by the test \texttt{test\_exponentials\_logarithms.c} in this folder.

\begin{enumerate}
\item \textbf{Purpose.} The file holds 64 true statements. Each says that an equation, an inequality or the domain of a function with exponentials or logarithms of the real variable $x$ holds exactly for the values printed as its result in a school textbook, for example $\frac{5^{x+1} - 1}{5^{x+1} - 5} < 0 \iff -1 < x < 0$.

\item \textbf{What it exercises.} The sign table with letters (rule 18p-8): an exponential $p^{x/L}$ or a logarithm $\log_b f(x)$ becomes a letter, the roots of the polynomials in the letter become comparisons in $x$, and a comparison has a value only where the arguments of its logarithms are positive. Also sums of logarithms turned into the logarithm of a product, monomials of exponentials in two bases, constants such as $\ln 3$ as endpoints, the logarithm of a power beside the logarithm of its base, the root of a product of exponentials, and a polynomial in the ratio of two exponentials.

\item \textbf{Variety.} Equations solved by substitution, by comparing exponents, with two bases; inequalities with fractions of exponentials; logarithmic equations with candidates outside the domain to discard; domains with square roots, nested logarithms and denominators. In exercise 691 the printed result is wrong: the statement here has the corrected one, and the printed one is in \texttt{false\_exponentials\_logarithms.tex}.

\item \textbf{Expected result.} The line \texttt{PASS: 64 of 64 statements proved, each with a written proof}.

\item \textbf{How to run it.} From the top folder of the repository:
\begin{verbatim}
cc -std=c11 -O2 -Wall -Wextra -o /tmp/test_exponentials_logarithms \
    tests/exponentials_logarithms/test_exponentials_logarithms.c \
    && /tmp/test_exponentials_logarithms
\end{verbatim}
\end{enumerate}
\end{document}
